\scanpage{682}

% 682-01, upper heading
$8^\circ)$ Return to the universal case, and the homomorphism
\[
 M^{\unsure{\natural|\sharp}}\longrightarrow
 M_{\rho,\sigma}^{\unsure{\natural|\sharp}}/\mu_\rho
\]

% 682-03, upper and middle text
We now suppose that $\unsure{\mathcal G}$ corresponds to the universal case
$\mathrm{GL}(2,\mathbb Z)^{\wedge}$, with an arbitrary
$\unsure{\mathcal G}^{\unsure{\natural|\sharp}}$. We shall denote the universal
case by a subscript $0$:
$\unsure{\mathcal G}_0,\ \unsure{\mathcal G}_0^+,\
\unsure{\mathcal G}_0^{\unsure{\natural|\sharp}}$ etc., and
$\unsure{\mathcal G}',\ \unsure{\mathcal G}'^{\unsure{\natural|\sharp}}$
etc. \unsure{the|the} \unsure{former} $\unsure{\mathcal G},\
\unsure{\mathcal G}^{\unsure{\natural|\sharp}}$ etc.

% 682-04, left margin at middle
\margin{$\unsure{\swarrow}\ Z_{\rho_0,\sigma_0}$ etc $\ldots$}
\margin{We shall write $\unsure{\mathcal G}_0^+$ instead of
$\unsure{\mathcal G}_0^{+\wedge}$ $\ldots$}

% 682-06, middle-lower text and display
I set
\[
\begin{array}{rl}
(107)\qquad M^{\unsure{\natural|\sharp}}&\subset M\\[-2pt]
\Vert&\\[-6pt]
\displaystyle
\left\{u\in M\ \middle|\
\begin{array}{ll}
\text{a) }u(\rho)\text{ is }\unsure{\mathcal G}^{\unsure{\natural|\sharp}}\text{-conjugate to}
 &\varepsilon_3[\tau](\rho)\\
\text{b) }u(\sigma)\quad\rule{4.2cm}{0.3pt}&\varepsilon_2[\tau](\sigma)\\
\text{c) }u(\varepsilon_0)\quad\rule{4.2cm}{0.3pt}&\varepsilon_0^\mu\\
\text{d) }u\text{ stabilizes }\unsure{M}^{\unsure{\natural|\sharp}}&
\end{array}\right\}&
\end{array}
\]
\margin{\uline{N.B.} $M^{\unsure{\natural|\sharp}}\unsure{\supset|\ni}\mu$}

% 682-09, lower text and display
where of course
\[
\mu=\chi(u),\quad
\varepsilon_2=\varepsilon_2(u)=\varepsilon_2(\mu),\quad
\varepsilon_3=\varepsilon_3(u)=\varepsilon_3(\mu).
\]
I say that $M^{\unsure{\natural|\sharp}}$ is a \uline{subgroup} of $M$.
We can see this by \unsure{going back over} the calculations on page\ $629$ and following.
I \unsure{prefer} to proceed as follows.

\margin{681}

% Printed archive footer.
\begin{center}
681\\
\texttt{https://grothendieck.umontpellier.fr}\\
Université de Montpellier, UMR CNRS 5149, Case courrier 051, Place Eugène Bataillon\\
34095 Montpellier cedex 5 - France
\end{center}


\scanpage{683}
\mspage{644}

% 683-02, upper text and display
We have seen that we have
\[
(108)\qquad
\underset{\deleted{\unsure{\sharp}}}{Z_{\rho_0,\sigma_0}}
\hookrightarrow M[0]\inserted{{}=M(0)\cdot L_0^{\unsure{\mathcal G}}}
\simeq \widetilde M_{0,3}[0]
=\widetilde M_{0,3}(0)\cdot L_0^{\unsure{\mathcal G}_{0,3}}
\]
an inclusion which is an equality if
$\unsure{\mathcal G}_0^{\unsure{\natural|\sharp}}=\unsure{\mathcal G}_0^+$,
and which identifies $Z_{\rho,\sigma}$ with
\[
M^{!}[0]\simeq M(0)\cdot L_0^\pi
\quad\text{if}\quad
\unsure{\mathcal G}_0^{\unsure{\natural|\sharp}}=\pi\unsure{:=}\Pi_0.
\]

% 683-08, bracketed passage
[ --- in fact, if
$\unsure{\mathcal G}_0^{\unsure{\natural|\sharp}}\subset\Pi_0$,
then we find
\[
\begin{split}
(109)\quad Z_{\rho,\sigma}=\left\{
u=\underset{\substack{\rotatebox{-90}{$\in$}\\M(0)}}{u_0}\cdot
\underset{\substack{\rotatebox{-90}{$\in$}\\\deleted{\unsure{\mathcal G}_0}L_0^\pi}}{\ell}
\ \middle|\
\begin{array}{l}
\text{a) }u(\rho)=\alpha(\rho)\quad\text{with }\alpha\in
\unsure{\mathcal G}_{\unsure{\bar\tau}}^{\unsure{\natural|\sharp}}\\
\text{b) }u(\sigma)=\beta(\sigma)\quad\text{with }\beta\in
\unsure{\mathcal G}_{\unsure{\bar\tau}}^{\unsure{\natural|\sharp}}\\
\text{d) }u\unsure{\varepsilon|\in}\text{ stabilizes }
\unsure{\mathcal G}_0^{\unsure{\natural|\sharp}}
\end{array}\right\}
\end{split}
\]
\margin{this explicit description does not \unsure{seem} to be of any use.}

% 683-12, middle-lower text with cancellations and additions
and we see that $\alpha,\beta$ in a) b) are determined by
$u=u_0\ell$ uniquely
(\deleted{as soon as} \unsure{even}
\unsure{if we|we} suppose only $u\in M^{!}[0]$ and no more,
and \inserted{\unsure{that we require}}
\deleted{\unsure{$\alpha\in$|with} $\unsure{\mathcal G}_0^{\unsure{\natural}}$}
$\alpha,\unsure{\beta}\in\unsure{\Pi_{0\bar\tau}}$,
\deleted{\unsure{can|In short}}
We know that this is a \unsure{subgroup|subgroup}
\inserted{of $\operatorname{Aut}(\unsure{\mathcal G}_0^+)\times
\widehat{\mathbb Z}^{*}\times\mu\times\mu$
--- and in fact, of $\operatorname{Aut}(\unsure{\mathcal G}_0^+)$}
(proposition\ page\ $629$).
\deleted{\unsure{Let us consider}}

% 683-19, lower paragraph
As this group operates on \deleted{$M$ and}
$\unsure{\mathcal G}$ via $\widetilde M_{0,3}$,
\inserted{$=\text{ via }\widetilde M_{0,3}^{!}[0]$}
it maps into $\widetilde M_{0,3}^{!}[0]$ by a homomorphism of groups,
hence into \deleted{$\unsure{\widetilde M_{0,3}[0]}$}
$M[0]^{!}\simeq\widetilde M_{0,3}[0]$ by a homomorphism of groups. ]

\margin{682}

% Printed archive footer.
\begin{center}
682\\
\texttt{https://grothendieck.umontpellier.fr}\\
Université de Montpellier, UMR CNRS 5149, Case courrier 051, Place Eugène Bataillon\\
34095 Montpellier cedex 5 - France
\end{center}


\scanpage{684}

% 684-01, upper text and tuple
\deleted{\unsure{The homomorphism}} Let us consider the system
\[
\begin{split}
(110)\quad\bigl(&\inserted{\unsure{\mathcal G}_0^{+\wedge},}
Z_{\rho_0,\sigma_0},\quad
L_0^{\unsure{\natural|\sharp}}\subset Z_{\rho_0,\sigma_0},\quad
Z_{\rho_0,\sigma_0}\longrightarrow
\operatorname{Aut}(\unsure{\mathcal G}_0^{\unsure{+\wedge}}),\quad
L_0^{\unsure{\natural|\sharp}}\hookrightarrow\unsure{\mathcal G}_0^{+\wedge}\bigr)
\end{split}
\]
(cf. ``digression'' below $9^\circ$) which gives rise to
\[
(110')\qquad
G_{\rho_0,\sigma_0},\quad
\unsure{\mathcal G}_0^{+\wedge}\subset G_{\rho_0,\sigma_0},\quad
Z_{\rho_0,\sigma_0}\hookrightarrow G_{\rho_0,\sigma_0},
\]
and let us consider the system analogous \inserted{to $(110^{*})$}
\[
\begin{split}
(111^{\unsure{*|\sharp}})\quad\bigl(&\inserted{\unsure{\mathcal G}_0^+,}
\underset{\substack{\Vert\\M(0)\cdot L_0^\pi}}{M[0]^{\unsure{!|\natural}}},\quad
\underset{\text{(\unsure{generated} by }\varepsilon_0\text{)}}{L_0^{\unsure{\mathcal G}}}
\subset M[0]^{\unsure{!|\natural}},\quad
\unsure{\deleted{M[0]}}^{\unsure{!|\natural}}\longrightarrow
\operatorname{Aut}(\unsure{\mathcal G}_0^+),\quad
L_0^{\unsure{\mathcal G}}\hookrightarrow\unsure{\mathcal G}_0^+\bigr)
\end{split}
\]
which gives rise (loc.\ cit ($9^\circ$)) to the situation
\[
(111')\qquad
\bigl(M,\quad\unsure{\mathcal G}_0^+\subset M,\quad
M[0]'\hookrightarrow\unsure{\mathcal G}_0^+\bigr).
\]

% 684-14, lower paragraph
Now thanks to $(108)$, we have a homomorphism from $(\unsure{110|109|119})$ into $(111)$,
corresponding to the identity on $\unsure{\mathcal G}_0^+$,
$(108)$ on $Z_{\rho_0,\sigma_0}$, the canonical inclusion\ :
\deleted{\unsure{on}}
$L_0^{\unsure{\natural|\sharp}}\subset L_0^{\unsure{\mathcal G}}$.
This yields a homomorphism from $(110')$ into $(111')$, in particular a canonical
homomorphism of exact sequences

\margin{683}

% Printed archive footer.
\begin{center}
683\\
\texttt{https://grothendieck.umontpellier.fr}\\
Université de Montpellier, UMR CNRS 5149, Case courrier 051, Place Eugène Bataillon\\
34095 Montpellier cedex 5 - France
\end{center}


\scanpage{685}
\mspage{645}

% 685-02, upper commutative diagram
% (112): the first \mathcal G_0^+ seems written over an earlier G_\rho (unresolved).
(112)\quad
\diagram{\begin{tikzcd}[column sep=2.1em,row sep=2.6em]
1 \arrow[r,no head] &
\unsure{\mathcal G}_0^+ \arrow[r] \arrow[d,"\unsure{1}"] &
G_{\rho_0,\sigma_0} \arrow[r] \arrow[d,hook] &
\gamma_{\rho_0,\sigma_0} \arrow[r] \arrow[d,hook,dotted] \arrow[d] & 1 \\
1 \arrow[r,no head] &
\unsure{\mathcal G}_0^+ \arrow[r,no head] & M \arrow[r] &
\widetilde\Gamma_{\mathbb Q} \arrow[r,no head] & 1
\end{tikzcd}}

% 685-03, explanatory diagram and text
where
\diagram{\begin{tikzcd}[column sep=4em,row sep=1em]
\gamma_{\rho_0,\sigma_0} \arrow[r] \arrow[d,phantom,"\unsure{\wr}" description] &
\widetilde\Gamma_{\mathbb Q} \arrow[d,phantom,"\unsure{\wr}" description] \\
Z_{\rho_0,\sigma_0}/L_0^{\unsure{\mu}} &
M[0]/L_0^{\unsure{\mathcal G}}
\end{tikzcd}}
is \inserted{\unsure{the one}} induced by
$(\unsure{108})$, and is identified with:
\[
(113)\qquad Z_{\rho_0,\sigma_0}/L_0^{\unsure{\natural|\sharp}}
\longrightarrow M[0]/L_0^{\unsure{\mathcal G}}
\]

% 685-07, middle text
By the very definition of
$L_0^{\unsure{\natural|\sharp}}$ as being
$Z_{\rho_0,\sigma_0}\cap\unsure{\mathcal G}_0^+$
\deleted{\unsure{and|this}}, we see that $(113)$ is injective\ hence
$G_{\rho_0,\sigma_0}\longrightarrow M$ is always injective
(it is an isomorphism\ if
$\unsure{\mathcal G}_0^{\unsure{\natural|\sharp}}\supset\Pi_0$
$\ldots$).
We deduce by comparison from $(112)$ a homomorphism of exact sequences

% 685-09, lower commutative diagram
(114)\quad
\diagram{\begin{tikzcd}[column sep=2.1em,row sep=2.6em]
1 \arrow[r,no head] &
\unsure{\mathcal G}_0^{\unsure{\natural|\sharp}} \arrow[r] \arrow[d,hook] &
G_{\rho_0,\sigma_0}^{\unsure{\natural|\sharp}} \arrow[r] \arrow[d,hook] &
\gamma_{\rho_0,\sigma_0} \arrow[r,no head] \arrow[d,hook] & 1 \\
1 \arrow[r,no head] &
\unsure{\mathcal G}_0^{\unsure{+}} \arrow[r] &
M \arrow[r] &
\deleted{\unsure{\gamma}}\widetilde\Gamma_{\mathbb Q} \arrow[r,no head] & 1.
\end{tikzcd}}

% 685-10, bottom text
We then see immediately that the image of
$G_{\rho_0,\sigma_0}^{\unsure{\natural|\sharp}}$ in $M$ is none other than
$M^{\unsure{\natural|\sharp}}$ $(107)$. We therefore have

\margin{684}

% Printed archive footer.
\begin{center}
684\\
\texttt{https://grothendieck.umontpellier.fr}\\
Université de Montpellier, UMR CNRS 5149, Case courrier 051, Place Eugène Bataillon\\
34095 Montpellier cedex 5 - France
\end{center}


\scanpage{686}

\[
(115)\qquad
\left\{
\begin{array}{ll}
\underset{\substack{\text{\unsure{defined in}}\\(107)}}{M^{\unsure{\natural|\sharp}}}
  \simeq G_{\rho_0,\sigma_0}^{\unsure{\natural|\sharp}},\\[5pt]
\gamma_{\rho_0,\sigma_0}\hookrightarrow\widetilde{\Gamma}_{\mathbb Q}
  \deleted{\unsure{\Leftarrow}\{u\in\widetilde{\Gamma}_{\mathbb Q}}
\end{array}
\right.
\]

\margin{
\diagram{
\begin{tikzcd}[row sep=small]
\text{\uline{to be checked}}\arrow[d]\\
\text{\uline{N.B.}\quad We have}
\end{tikzcd}
}
\[
(116)\quad
\left\{
\begin{array}{l}
\gamma_{\rho_0,\sigma_0}
 \deleted{\unsure{=\Gamma_0}}
 \inserted{\unsure{=\Gamma_0^{\natural}|\Gamma_0^{\sharp}}}
 =\left\{u\in\widetilde{\Gamma}_{\mathbb Q}\ \middle|\
 \begin{array}{l}
\text{the lifting }u_0\\
\text{of }u\text{ to }u_0\in M(0)\\
\text{satisfies the}\\
\text{conditions a) b) d)}\\
\text{of (107)}
\end{array}\right\}\\
\text{(\unsure{notations}}\\
\gamma_0=\Gamma_0^{\unsure{\natural|\sharp}},
\ \text{where }\Gamma_0\text{ denotes }\widetilde{\Gamma}_{\mathbb Q}\text{)}
\end{array}
\right.
\]
}

which \unsure{shows} at the same time in the
\unsure{process} that $M^{\unsure{\natural|\sharp}}$ is indeed a
subgroup of $M$.

We now have
\[
(\unsure{117})\qquad
M_{\rho_0,\sigma_0}^{\unsure{\natural|\sharp}}
\overset{\unsure{\mathrm{dfn}|\mathrm{déf}}}{\simeq}
\unsure{N|M}_{\rho_0}^{\unsure{\natural|\sharp}}
\times_{\gamma_{\rho_0,\sigma_0}}
N_{\sigma_0}^{\unsure{\natural|\sharp}}
\times_{\unsure{\gamma_{\rho_0,\sigma_0}}}
\underbrace{G_{\rho_0,\sigma_0}^{\unsure{\natural|\sharp}}}
 _{\simeq M^{\unsure{\natural|\sharp}}},
\]
where $G_{\rho_0,\sigma_0}^{\unsure{\natural|\sharp}}\simeq
M^{\unsure{\natural|\sharp}}$ \unsure{therefore|gives}
\[
\begin{aligned}
N_{\rho_0}^{\unsure{\natural|\sharp}}
 &\unsure{\simeq|=}\left\{u\in M^{\unsure{\natural|\sharp}}\mid
 u(\rho)=\rho^{\varepsilon_3(u)}\right\}\\
N_{\sigma_0}^{\unsure{\natural|\sharp}}
 &\unsure{\simeq|=}\left\{u\in M^{\unsure{\natural|\sharp}}\mid
 u(\sigma)\unsure{\equiv|=}\sigma^{\varepsilon_2(u)}\right\}
\end{aligned}
\]
and finally we have
\[
(118)\qquad
\left\{
\begin{array}{l}
M_{\rho_0,\sigma_0}^{\unsure{\natural|\sharp}}
\simeq
N_{M^{\unsure{\natural|\sharp}}}(\rho)
\times_{\deleted{\unsure{\gamma_{\rho_0,\sigma_0}}}
 \inserted{\Gamma_0^{\unsure{\natural|\sharp}}}}
N_{M^{\unsure{\natural|\sharp}}}(\sigma)
\times_{\deleted{\unsure{\gamma_{\rho_0,\sigma_0}}}
 \inserted{\Gamma_0^{\unsure{\natural|\sharp}}}}
M^{\unsure{\natural|\sharp}}.\\[5pt]
\text{where }\Gamma_0^{\unsure{\natural|\sharp}}\unsure{\ne|=}\gamma_0
\simeq\gamma_{\rho_0,\sigma_0}\subset\widetilde{\Gamma}_{\mathbb Q}
\quad\text{\unsure{made explicit} in (116)},
\end{array}
\right.
\]
\uuline{whence} an \deleted{isomorphism of} \unsure{exact|exact}
\unsure{sequences|sequence}
\[
(119)
\]
\diagram{
\begin{tikzcd}[column sep=small]
1\arrow[r]&
\overbrace{SZ_{\rho_0}\times SZ_{\sigma_0}}^{\inserted{L_\rho\times L_\sigma}}
\arrow[r]&
M_{\rho_0,\sigma_0}^{\unsure{\natural|\sharp}}
\arrow[r]\arrow[d,"\inserted{\unsure{\natural|\sharp}}"']&
\deleted{\unsure{\gamma}}M^{\unsure{\natural|\sharp}}\arrow[r]&1\\
&&{}&&
\end{tikzcd}
}
where the kernel $SZ_{\rho_0}\times SZ_{\sigma_0}$ is a
\unsure{very small} group, namely

\margin{685}

% Printed archive footer.
\begin{center}
685\\
\texttt{https://grothendieck.umontpellier.fr}\\
Université de Montpellier, UMR CNRS 5149, Case courrier 051, Place Eugène Bataillon\\
34095 Montpellier cedex 5 - France
\end{center}